% ── Appendix: Technical reproducibility and extended diagnostics ─────────────
\section{Technical Reproducibility and Extended Diagnostics}
\label{sec:appendix_repro}

This appendix records the technical details needed to interpret and reproduce the report's main
results. The accompanying GitHub repository \citep{faverio2026coastal_gap_repository} contains the
complete scripts, feature matrices, per-gap prediction tables, model outputs, and extended diagnostic
files. The appendix instead reports the target definitions, validation supports, product provenance,
benchmark backing tables, and selected robustness diagnostics needed to make the main text auditable.

\subsection{Target definitions and daily quality control}
\label{app:targets_qc}

Both targets are daily means computed from hourly Tongoy buoy (BTG) observations from the CEAZAMET
BTG station \citep{ceazamet_btg_station}. A day is eligible only if at least 18 hourly observations
are valid. Chlorophyll is modelled and scored on $\log_{10}(\texttt{chl\_mean})$, while dissolved
oxygen is modelled and scored directly in mg/L.

\begin{table}[h]
  \centering
  \small
  \caption{Daily target definitions and quality-control summaries. The date span for both targets
  is 2015-07-01 to 2026-05-31, for a total of 3988 possible daily records.}
  \label{tab:appendix_target_qc}
  \begin{tabularx}{\textwidth}{@{}l l l r r r r@{}}
    \toprule
    Target & Channel & Unit & Eligible days & Missing days & Missing share & Real gaps \\
    \midrule
    Chlorophyll-a
    & BTGCLF
    & mg~m$^{-3}$
    & 3012
    & 976
    & 24.5\%
    & 128 \\

    Dissolved oxygen
    & BTGOXD2
    & mg/L
    & 2880
    & 1108
    & 27.8\%
    & 125 \\
    \bottomrule
  \end{tabularx}
\end{table}

\begin{table}[h]
  \centering
  \small
  \caption{Distribution summaries for eligible daily target values.}
  \label{tab:appendix_target_distributions}
  \begin{tabular}{l r r r r}
    \toprule
    Target & Median & IQR & p05 & p95 \\
    \midrule
    Chlorophyll-a (mg~m$^{-3}$)
    & 3.48 & 5.04 & 0.79 & 15.93 \\

    Dissolved oxygen (mg/L)
    & 6.43 & 2.34 & 2.82 & 8.80 \\
    \bottomrule
  \end{tabular}
\end{table}

\subsection{External products and provenance}
\label{app:data_products}

External products are used as predictor channels. Table~\ref{tab:appendix_products}
extends the compact product summary in the main text.

\begin{table}[h]
  \centering
  \small
  \caption{External products used to construct report predictor channels.}
  \label{tab:appendix_products}
  \begin{tabularx}{\textwidth}{@{}l X l l X@{}}
    \toprule
    Product / source & Role in report & Native scale & Coverage & Notes \\
    \midrule
    CEAZAMET BTG / PLV
    & In-situ sensor targets and local atmospheric forcing.
    & point, hourly
    & --
    & BTG provides the chlorophyll and oxygen targets; PLV provides local meteorological forcing. \\

    GlobColour L4 \chl{}
    & Satellite biological proxy for surface chlorophyll.
    & $\sim$4~km, daily
    & $\sim$90\%
    & \\

    MUR SST
    & Sea-surface temperature and thermal-state predictors.
    & $\sim$1~km, daily
    & $\sim$98\%
    & Used for SST and thermal forcing summaries. \\

    CMEMS winds
    & Regional wind and upwelling forcing predictors.
    & $\sim$25~km, hourly
    & $\sim$90\%
    & Used for wind speed, stress and upwelling-related summaries. \\

    GLORYS12 currents
    & Regional current and transport predictors.
    & $\sim$8~km, daily
    & $>99\%$
    & Used for current speed, components and transport summaries. \\
    \bottomrule
  \end{tabularx}
\end{table}

\subsection{Artificial-gap pools and validation supports}
\label{app:gap_pools}

\begin{table}[h]
  \centering
  \small
  \caption{Artificial-gap pools and report validation supports.}
  \label{tab:appendix_gap_supports}
  \begin{tabularx}{\textwidth}{@{}l l l r X@{}}
    \toprule
    Target & Support & Gap lengths & Gaps & Role in report \\
    \midrule
    Chlorophyll
    & Full TS-ICL pool
    & $1,3,7,10,14,21,30,45,60$
    & 681
    & TS-ICL calibration and covariate-signal diagnostics. \\

    Chlorophyll
    & Shared ladder support
    & $1,3,7,14,30$
    & 449
    & Main chlorophyll method ranking; all headline comparators evaluated on the same hidden intervals. \\

    Oxygen
    & Primary support
    & $1,3,7,10,14,21,30$
    & 406
    & Main oxygen method ranking and distribution-tail diagnostics. \\

    Oxygen
    & Exploratory longer support
    & $>30$
    & 6
    & Not used for headline oxygen claims. \\
    \bottomrule
  \end{tabularx}
\end{table}
\vspace{-0.5cm}
\subsection{Feature construction}
\label{app:feature_model_details}

The main text groups predictors into scientific information channels: calendar/seasonality, target
history, satellite biological proxy, wind/upwelling forcing, SST, and
current/transport information. Tabular models require these channels to be converted into explicit
per-day feature vectors. TS-ICL accepts time-indexed covariate channels, which may still be engineered summaries, but they enter as channels over time rather than as one fixed row of
tabular features per hidden day.

\begin{small}
\setlength{\tabcolsep}{4pt}
\renewcommand{\arraystretch}{1.18}

\begin{longtable}{P{2.4cm}P{3.2cm}P{4.9cm}P{4.4cm}P{1cm}}  \caption{Feature-construction inventory.}
  \label{tab:appendix_feature_construction} \\
  \toprule
  Feature family & Raw source & Constructed features & Physical meaning & Target \\
  \midrule
  \endfirsthead

  \toprule
  Feature family & Raw source & Constructed features & Physical meaning & Target \\
  \midrule
  \endhead

  \bottomrule
  \endfoot

  Calendar / seasonality
  & Timestamp
  & \begin{tabitemize}[leftmargin=*, nosep, topsep=0pt]
    \item month
    \item year
    \item sine/cosine doy
  \end{tabitemize}
  & Position within the annual upwelling and productivity cycle.
  & Both \\

  Target history
& Target record before the gap.
& \begin{tabitemize}[leftmargin=*, nosep, topsep=0pt]
    \item 3-, 7-, and 14-day pre-gap statistics;
    \item short pre-gap slopes;
  \end{tabitemize}
& Trend of the observed target approaching the missing interval.
& Both \\

  Gap-edge context
  & Target record on both gap sides
  & \begin{tabitemize}[leftmargin=*, nosep, topsep=0pt]
    \item first pre/post-gap values;
    \item gap-position geometry;
    \item edge-to-edge mean, slope;
  \end{tabitemize}
  & Behaviour bracketing an internal gap, used to correct interpolation.
  & Both \\

  Satellite biological proxy
  & GlobColour L4 \chl{} \citep{garnesson2019globcolour}.
  & \begin{tabitemize}[leftmargin=*, nosep, topsep=0pt]
    \item nearest-cell log-\chl{};
    \item $3\times3$ / $5\times5$ box summaries;
    \item patchiness, gradients;
    \item doy/moy anomalies;
  \end{tabitemize}
  & Regional surface chlorophyll as an external proxy for the in-situ biological signal.
  & CLF \\

  Wind / upwelling forcing
  & PLV station winds and CMEMS winds \citep{copernicus_wind_l4}
  & \begin{tabitemize}[leftmargin=*, nosep, topsep=0pt]
    \item wind components/speed;
    \item upwelling indices;
    \item wind-direction persistence;
    \end{tabitemize}
  & Mechanical forcing of coastal upwelling.
  & Both \\

  Solar / atmospheric forcing
  & PLV meteorology.
  & \begin{tabitemize}[leftmargin=*, nosep, topsep=0pt]
    \item radiation lags and rolling means;
    \item pressure/humidity summaries;
  \end{tabitemize}
  & Local radiative and atmospheric state over the bay.
  & Both \\

  SST / thermal state
  & MUR SST \citep{jpl_mur_sst}
  & \begin{tabitemize}[leftmargin=*, nosep, topsep=0pt]
    \item nearest-cell SST;
    \item cooling indices;
    \item coastal gradients;
  \end{tabitemize}
  & Thermal tracer of recently upwelled water and surface stratification.
  & Both \\

  Current / transport
  & GLORYS12 currents \citep{lellouche2021glorys12,copernicus_glorys12}
  & \begin{tabitemize}[leftmargin=*, nosep, topsep=0pt]
    \item current speed/components;
    \item alongshore and cross-shore;
    \item divergence, vorticity, Okubo Weiss;
  \end{tabitemize}
  & Advection, exchange, and retention conditions over the shelf.
  & Both \\

  Local BTG diagnostics
  & BTG water temperature and pressure.
  & \begin{tabitemize}[leftmargin=*, nosep, topsep=0pt]
    \item daily mean;
    \item valid hour count;
  \end{tabitemize}
  & Co-located physical state at the buoy, distinct from the oxygen sensor.
  & OXY \\
\end{longtable}
\end{small}

\subsection{Chlorophyll benchmark backing tables}
\label{app:chl_benchmark_tables}

This subsection gives the numerical backing for the matched chlorophyll benchmark in
Section~\ref{sec:chl_benchmark}. All values are computed on $\log_{10}(\texttt{chl\_mean})$.

\begin{table}[h]
  \centering
  \small
  \caption{Chlorophyll benchmark summary.}
  \label{tab:appendix_chl_matched_metrics}
  \begin{tabular}{l r}
    \toprule
    Method & MAE \\
    \midrule
    TS-ICL satellite-proxy & 0.221 \\
    TS-ICL target-only & 0.225 \\
    Gap-edge residual & 0.229 \\
    Gaussian process & 0.234 \\
    Linear interpolation & 0.239 \\
    External tabular (ExtraTrees) & 0.277 \\
    External tabular (HGB) & 0.287 \\
    \bottomrule
  \end{tabular}
\end{table}

\begin{table}[h]
  \centering
  \small
  \caption{Chlorophyll MAE by gap length.}
  \label{tab:appendix_chl_by_length}
  \begin{tabular}{r r r r r r r}
    \toprule
    $L$ & Linear interp. & Gaussian process & External tabular & Gap-edge & TS target & TS satellite \\
    \midrule
    1  & 0.108 & 0.104 & 0.248 & 0.127 & 0.106 & 0.106 \\
    3  & 0.133 & 0.129 & 0.257 & 0.143 & 0.123 & 0.122 \\
    7  & 0.190 & 0.199 & 0.274 & 0.185 & 0.186 & 0.180 \\
    14 & 0.254 & 0.248 & 0.267 & 0.239 & 0.239 & 0.235 \\
    30 & 0.278 & 0.266 & 0.295 & 0.265 & 0.259 & 0.254 \\
    \bottomrule
  \end{tabular}
\end{table}

\begin{table}[H]
  \centering
  \small
  \caption{Paired chlorophyll MAE differences against TS-ICL satellite-proxy.}
  \label{tab:appendix_chl_paired_deltas}
  \begin{tabular}{l r r}
    \toprule
    Comparator & $\Delta$MAE & 95\% CI \\
    \midrule
    Linear interpolation & -0.018 & [-0.028,\,-0.008] \\
    External tabular (ExtraTrees) & -0.057 & [-0.072,\,-0.043] \\
    Gaussian process & -0.013 & [-0.020,\,-0.006] \\
    Gap-edge residual & -0.009 & [-0.016,\,-0.001] \\
    TS-ICL target-only & -0.005 & [-0.007,\,-0.002] \\
    \bottomrule
  \end{tabular}
\end{table}

\subsection{Chlorophyll event diagnostics}
\label{app:chl_extended_diagnostics}

High-\chl{} events remain difficult for all methods. Table~\ref{tab:appendix_chl_event_mae}
reports the event/non-event split used to support Figure~\ref{fig:chl_event_limitation}.

\begin{table}[H]
  \centering
  \small
  \caption{Chlorophyll event and non-event MAE.}
  \label{tab:appendix_chl_event_mae}
  \begin{tabular}{l r r}
    \toprule
    Method & Non-event MAE & High-\chl{} event MAE \\
    \midrule
    Linear interpolation & 0.223 & 0.272 \\
    Gaussian process & 0.204 & 0.295 \\
    External tabular (ExtraTrees) & 0.235 & 0.366 \\
    Gap-edge residual & 0.205 & 0.281 \\
    TS-ICL target-only & 0.202 & 0.274 \\
    TS-ICL satellite-proxy & 0.199 & 0.266 \\
    \bottomrule
  \end{tabular}
\end{table}

\subsection{Oxygen benchmark backing tables}
\label{app:oxygen_benchmark_tables}

This subsection gives the numerical backing for the primary oxygen benchmark in
Section~\ref{sec:oxygen_benchmark}. All oxygen MAE values are gap-weighted and reported in mg/L.

\begin{table}[H]
  \centering
  \small
  \caption{Oxygen benchmark by gap length.}
  \label{tab:appendix_oxygen_by_length}
  \begin{tabular}{r r r r r r r r}
    \toprule
    $L$ & $n$ & Interp. & Ext. tabular & GP & Gap-edge & TS target & TS physical \\
    \midrule
    1  & 100 & 0.339 & 0.606 & 0.345 & 0.380 & 0.307 & 0.300 \\
    3  & 100 & 0.635 & 0.805 & 0.642 & 0.674 & 0.618 & 0.562 \\
    7  & 78  & 0.863 & 0.949 & 0.861 & 0.876 & 0.846 & 0.824 \\
    10 & 55  & 0.970 & 0.962 & 0.970 & 0.971 & 0.949 & 0.930 \\
    14 & 36  & 1.078 & 0.998 & 1.086 & 1.094 & 1.027 & 0.949 \\
    21 & 22  & 0.966 & 1.102 & 0.963 & 1.005 & 0.992 & 0.945 \\
    30 & 15  & 1.312 & 1.040 & 1.346 & 1.306 & 1.171 & 1.153 \\
    \bottomrule
  \end{tabular}
\end{table}

\begin{table}[h]
  \centering
  \small
  \caption{Paired oxygen MAE differences against TS-ICL physical covariates.}
  \label{tab:appendix_oxygen_paired_deltas}
  \begin{tabular}{l r r}
    \toprule
    Comparator & $\Delta$MAE (mg/L) & 95\% CI \\
    \midrule
    Linear interpolation & -0.059 & [-0.086,\,-0.032] \\
    External tabular (ExtraTrees) & -0.172 & [-0.229,\,-0.121] \\
    Gaussian process & -0.064 & [-0.091,\,-0.036] \\
    Gap-edge residual (ExtraTrees) & -0.084 & [-0.114,\,-0.056] \\
    TS-ICL target-only & -0.032 & [-0.051,\,-0.014] \\
    \bottomrule
  \end{tabular}
\end{table}

The external-tabular curve appears favourable at $L=30$ in the by-length plot, but this stratum has
only 15 gaps. A focused robustness audit was therefore used to decide whether this local reversal
should be treated as a headline result. Table~\ref{tab:appendix_oxygen_l30_audit} shows why it is
kept as a diagnostic caveat only.

\begin{table}[h]
  \centering
  \small
  \caption{Oxygen $L=30$ external-tabular robustness audit.}
  \label{tab:appendix_oxygen_l30_audit}
  \begin{tabularx}{\textwidth}{@{}l X@{}}
    \toprule
    Diagnostic & Result \\
    \midrule
    Number of $L=30$ gaps
    & 15 \\

    Bootstrap 95\% CI for mean paired difference
    & [-0.644,\,+0.092] mg/L; includes zero. \\

    Most influential gap
    & \texttt{OX\_L030\_20241115}, where interpolation has very high error. \\

    Sensitivity interpretation
    & The apparent mean reversal is influenced by a small number of high-error interpolation
      cases. \\

    Median / trimmed comparison
    & Alternative summaries still show some external-tabular advantage, but the mean advantage is
      not confidence-interval resolved. \\

    Report interpretation
    & Small-stratum diagnostic, not a headline reversal of the pooled conclusion. \\
    \bottomrule
  \end{tabularx}
\end{table}

\subsection{Oxygen tail diagnostics}
\label{app:oxygen_extended_diagnostics}

The oxygen distribution-tail diagnostics use empirical quantile bands, not ecological thresholds.
Relative improvement is computed against interpolation.
The high-tail loss in Table~\ref{tab:appendix_oxygen_quantile_bands} is confidence-interval
resolved, while the low-tail loss is a point-estimate loss whose confidence interval includes zero.
The tail-persistence split sharpens the diagnostic pattern but uses smaller strata, especially for
sustained high oxygen, and is therefore not treated as a separate headline comparison.

\begin{table}[h]
  \centering
  \small
  \caption{Oxygen TS-ICL performance by empirical oxygen quantile band.}
  \label{tab:appendix_oxygen_quantile_bands}
  \begin{tabular}{l r r r}
    \toprule
    Band & $n$ & Relative improvement & 95\% CI \\
    \midrule
    $<p10$ ($<3.78$ mg/L) & 123 & -9.8\% & [-25.7,\,+2.2]\% \\
    $p10$--$p25$ (3.78--5.10 mg/L) & 170 & +25.5\% & [+19.0,\,+32.2]\% \\
    $p25$--$p50$ (5.10--6.43 mg/L) & 223 & +21.6\% & [+15.3,\,+27.6]\% \\
    $p50$--$p75$ (6.43--7.44 mg/L) & 211 & +6.7\% & [-0.8,\,+13.9]\% \\
    $p75$--$p90$ (7.44--8.30 mg/L) & 121 & -5.9\% & [-20.3,\,+6.7]\% \\
    $>p90$ ($>8.30$ mg/L) & 79 & -18.8\% & [-35.8,\,-5.6]\% \\
    \bottomrule
  \end{tabular}
\end{table}

\begin{table}[H]
  \centering
  \small
  \caption{Oxygen tail-persistence diagnostics.}
  \label{tab:appendix_oxygen_tail_persistence}
  \begin{tabular}{l r r}
    \toprule
    Tail-persistence stratum & $n$ & Relative improvement vs interpolation \\
    \midrule
    Low tail, isolated & 25 & 23.1\% \\
    Low tail, short run (2--6 d) & 76 & -5.2\% \\
    Low tail, sustained run ($\geq$7 d) & 50 & -22.3\% \\
    High tail, isolated & 32 & -12.3\% \\
    High tail, short run (2--6 d) & 45 & -13.8\% \\
    High tail, sustained run ($\geq$7 d) & 12 & -123.4\% \\
    \bottomrule
  \end{tabular}
\end{table}